% ============================================================
% GDP 논문 (한국어) — IEEE CASE 2026
% 피드백 반영 수정본
% Overleaf 복붙용 LaTeX 텍스트
% ============================================================

% ------------------------------------------------------------
\begin{abstract}
% ------------------------------------------------------------
스마트 인더스트리(Smart Industry) 시대의 핵심 기술인 디지털 트윈은 정교한 제조 데이터 확보와 시뮬레이션 프레임워크를 필요로 한다. 그러나 실제 산업 데이터는 기업의 핵심 보안 자산으로서 외부 접근이 제한되어 있고, 설계 정보를 분석 도구에 연결하는 인터페이스 비용이 과도하게 높다.

본 논문에서는 대규모 언어 모델(LLM)을 활용하여 공정 설계의 진입 장벽을 낮추고 도구 간 연동을 자동화하는 프레임워크인 GDP(Generative Digital-twin Prototyper)를 제안한다. GDP는 Siemens가 제안한 BOP(Bill of Process) 개념을 참조하되 LLM 생성 효율을 극대화한 독자적 공정 데이터 모델을 채택하며, 두 가지 사용자 그룹을 지원한다: (1)~제조 지식이 부족한 \textit{비전문가}가 제로샷 방식으로 연구용 모의 생산 라인을 설계할 수 있고, (2)~\textit{전문가}가 기존 스크립트 어댑터 자동 생성, AI 기반 분석 도구 합성, 스키마 가이드 최적화의 세 가지 도구 연동 방식을 통해 인터페이스 개발 비용을 대폭 절감할 수 있다. 특히 LLM이 실행 오류를 자율적으로 수정하는 Auto-Repair 메커니즘을 통해 연동 강건성을 확보한다.

실험 결과, 5종의 최신 LLM을 사용한 제로샷 공정 생성에서 최고 모델(Gemini 2.5 Flash)이 81.2\%의 정확도를 달성하였다. 전문가 인터뷰를 통해 도구 연동 기능의 실용적 가치를 확인하였으며, 비전문가가 생성한 모의 공정 데이터가 휴머노이드 로봇 및 피지컬 AI(Physical AI) 학습을 위한 가상 테스트베드로 활용될 잠재성을 제시한다.
\end{abstract}

\begin{IEEEkeywords}
Digital Twin, LLM, Smart Industry, Bill of Process, Tool Integration, Auto-Repair, Physical AI
\end{IEEEkeywords}

%% ============================================================
\section{서론}
\label{sec:introduction}
%% ============================================================

스마트 인더스트리(Smart Industry) 시대의 제조 경쟁력은 물리적 공장을 가상 공간에 복제하여 사전 검증하는 디지털 트윈 기술의 성숙도에 좌우된다. 디지털 트윈은 설계 단계에서의 병목 예측, 공정 최적화, 그리고 실시간 모니터링을 가능하게 하여 제조 혁신의 핵심 인프라로 자리매김하고 있다~\cite{b1}. 그러나 연구 개발 단계에서 디지털 트윈을 진전시키기 위해서는 두 가지 주요 장벽을 극복해야 한다.

첫째는 \textbf{벤치마크 데이터의 부재와 파편화}이다. 실제 제조 데이터는 기업의 핵심 보안 자산으로 공개가 엄격히 제한되어 있어, 외부 연구자와 개발자가 시뮬레이션 발전을 위한 신뢰성 있는 공개 데이터셋을 확보하기 어렵다. 예를 들어, 반도체 후공정이나 배터리 셀 제조와 같은 첨단 산업의 공정 데이터는 기술 유출 방지를 위해 거의 공개되지 않으며, 공개된 데이터조차 형식과 세분화 수준이 상이하여 직접적인 시뮬레이션 입력으로 활용하기 어렵다. 이로 인해 디지털 트윈 연구는 소수의 단순화된 벤치마크에 의존하거나, 연구자 개인이 수작업으로 모의 데이터를 구축해야 하는 비효율적 상황에 처해 있다.

둘째는 공정 설계와 분석 도구 간의 \textbf{인터페이스 파편화}이다. 제조 현장에서는 라인 밸런싱, 병목 분석, 에너지 추정, 레이아웃 최적화 등 다양한 분석 도구가 활용되지만, 각 도구는 고유한 입출력 형식을 요구한다. 설계 데이터가 이러한 시뮬레이션이나 최적화 엔진으로 유입되기 위해서는 복잡한 데이터 변환이 필수적이며, 이로 인해 엔지니어들은 핵심 분석 업무가 아닌 데이터 전처리 및 어댑터 개발에 과도한 시간을 소비한다. 실제로 도구 하나를 새로운 데이터 모델에 연동하는 데 수일의 개발 기간이 소요되는 경우가 빈번하다.

이러한 장벽은 두 가지 상이한 사용자 그룹에 각기 다른 방식으로 영향을 미친다.

\textit{비전문가}---로보틱스, 컴퓨터 비전 등 타 분야 연구자나 교육 목적 사용자---는 제조 도메인 지식의 부족으로 인해 시뮬레이션이나 디지털 트윈 연구에 필요한 모의 공정 데이터 자체를 생성하지 못한다. 이들은 제조 공정의 세부 단계, 설비 구성, 자재 흐름 등에 대한 전문 지식이 없으므로 현실적인 모의 생산 라인을 구축하는 것이 사실상 불가능하다.

\textit{전문가}---공정 엔지니어, 시뮬레이션 연구자 등---는 공정 지식은 보유하고 있으나, 자신이 개발하거나 보유한 분석 도구를 새로운 데이터 모델에 연동하는 반복적 인터페이스 개발에 과도한 노력을 소비한다. 도구마다 입출력 스키마가 상이하고, 데이터 변환 과정에서 발생하는 런타임 오류의 디버깅이 추가적 부담으로 작용한다.

본 연구는 이 두 가지 문제를 동시에 해결하는 GDP(Generative Digital-twin Prototyper) 프레임워크를 제안한다. GDP는 LLM 최적화 공정 모델을 인터페이스 브릿지로 활용하여, 비전문가에게는 자연어 기반 제로샷 공정 설계를, 전문가에게는 자동화된 도구 연동을 제공한다. 본 연구의 주요 기여는 다음과 같다:

\begin{enumerate}
\item \textit{시뮬레이션 적합 데이터 모델 정제}: Siemens가 제안한 BOP(Bill of Process) 개념을 참조하되, 복잡한 산업 구조에서 시뮬레이션에 필수적인 속성만을 추출하여 LLM 기반 생성과 분석 도구 연동에 적합한 경량 데이터 구조를 확립한다.

\item \textit{전문가를 위한 인터페이스 자동화 및 Auto-Repair}: 기존 스크립트 어댑터 자동 생성, AI 기반 분석 도구 합성, 스키마 가이드 최적화의 세 가지 연동 방식을 제공하고, LLM이 실행 시 런타임 오류를 자율적으로 감지·수정하는 자가 수리 메커니즘을 구현한다.

\item \textit{비전문가를 위한 제로샷 공정 설계 및 피지컬 AI 데이터 생성}: 제조 지식이 부족한 사용자도 자연어 입력만으로 구조적으로 완결된 BOP를 생성할 수 있으며, 이를 통해 생성된 모의 데이터는 피지컬 AI 학습을 위한 가상 제조 환경 구축에 활용될 수 있다.
\end{enumerate}

특히, 비전문가가 GDP를 통해 생성한 다양한 모의 생산 라인 데이터는 휴머노이드 로봇 및 피지컬 AI(Physical AI) 시스템이 가상 환경에서 제조 작업을 사전 학습하는 테스트베드로 활용될 잠재성을 지닌다. 최근 NVIDIA의 Omniverse, Google DeepMind의 로봇 시뮬레이션 등에서 볼 수 있듯이, 가상 제조 환경에서의 대규모 학습 데이터 생성은 피지컬 AI 발전의 핵심 과제로 부상하고 있다. GDP는 실제 공장 데이터 없이도 다양한 제조 시나리오를 대규모로 생성하여, 이러한 AI 학습 데이터의 다양성과 규모를 확보하는 데 기여할 수 있다.

본 논문의 나머지 구성은 다음과 같다. II장에서 관련 연구를 검토하고, III장에서 GDP 프레임워크의 아키텍처를 제시한다. IV장에서 실험 결과를 분석하며, V장에서 결론과 향후 연구 방향을 논의한다.

%% ============================================================
\section{관련 연구}
\label{sec:related}
%% ============================================================

\subsection{디지털 트윈 표준화 및 모델링 자동화}

디지털 트윈은 물리적 자산, 프로세스 또는 시스템의 가상 복제본으로, 모니터링과 시뮬레이션을 통한 최적 의사결정을 지원한다~\cite{b1}. 초기 항공우주 분야 중심의 연구에서 최근 제조 가시성 향상을 위한 ``디지털 트윈 Shop-floor'' 개념으로 확장되었으며~\cite{b2}, ISO 23247 표준은 제조 디지털 트윈을 위한 4계층 아키텍처를 제시하여 데이터 수집과 가상 모델 간 연결성을 강조하였다~\cite{b3}.

그러나 이러한 표준 모델을 실제 현장에 적용하는 데에는 여전히 높은 모델링 비용이 수반된다. Uhlemann 등은 물리적 설비 변경을 가상 모델에 실시간 반영하는 공정 모델링 단계가 여전히 숙련 엔지니어의 수작업에 의존함을 지적하였다~\cite{b4}. 기존 온톨로지 기반 접근법은 제조 지식의 형식화를 시도하였으나, 사전 정의된 규칙의 한계를 넘어 복잡한 신규 공정 시나리오를 유연하게 수용하기 어렵다~\cite{b5}.

한편, Siemens는 제조 공정을 체계적으로 기술하기 위한 BOP(Bill of Process) 개념을 제안하였다~\cite{b7}. BOP는 Plant--Line--Station--Operation의 계층 구조를 통해 제조 공정의 순서, 리소스 배정, 시간 정보를 구조화하며, Tecnomatix 등의 시뮬레이션 환경에서 공정 모델링의 기반 데이터로 활용된다. 그러나 이 심층 계층 구조는 수작업 입력 비용이 높고, 외부 도구와의 연동 시 복잡한 데이터 매핑을 요구한다는 한계가 있다. 본 연구는 이 BOP 개념을 참조하되, LLM 기반 자동 생성에 최적화된 평탄화 구조로 정제하여 활용한다.

\subsection{LLM 기반 제조 지식 구조화}

GPT-4 등의 대규모 언어 모델은 비정형 텍스트에서 논리적 구조를 추출하여 산업 데이터로 변환하는 탁월한 능력을 입증하였다~\cite{b6}. 제조 분야에서는 LLM을 활용한 컴퓨터 지원 공정 계획(CAPP) 연구가 제조 매뉴얼이나 설계 도면으로부터 작업 시퀀스를 자동 추출할 가능성을 열었다. 또한 LLM의 제로샷 추론 능력을 활용하여 복잡한 제조 시나리오를 JSON 또는 XML 형식의 구조화된 공정 명세로 인간 개입 없이 변환하는 연구가 진행되고 있다~\cite{b8}.

특히, 최근 LLM의 구조화된 출력(structured output) 생성 능력이 크게 향상되면서, 테이블 형식 데이터의 이해~\cite{b8}와 스키마 준수 JSON 생성이 실용적 수준에 도달하였다~\cite{b9}. 이는 LLM이 단순한 텍스트 생성을 넘어, 공정 데이터 모델과 같은 복잡한 구조적 제약 조건을 만족하는 출력을 생성할 수 있음을 시사한다.

그러나 단순 텍스트 추출을 넘어 Siemens가 제안한 BOP와 같은 공정 데이터 구조에 정밀하게 매핑하여 시뮬레이션 도구와 즉시 연동 가능한 출력을 생성하는 연구는 초기 단계에 머물러 있다. 대부분의 기존 연구는 LLM 출력의 형식적 정확성에 초점을 맞추고 있으며, 생성된 공정 데이터가 후속 분석 도구에 실제로 연동 가능한지를 검증하는 엔드-투-엔드 파이프라인은 거의 제안되지 않았다. 본 연구는 LLM의 사전 지식(prior knowledge)을 공정 모델 정제 파이프라인과 결합하여 이 격차를 해소하는 접근법을 제시한다.

\subsection{이기종 도구 연동 및 자동 어댑터 생성}

공정 설계 데이터가 시뮬레이션(예: Plant Simulation)이나 최적화 엔진으로 전달될 때 발생하는 ``인터페이스 파편화''는 디지털 트윈 구축의 최대 병목이다~\cite{b10}. 전통적으로 AutomationML이나 OPC UA 등의 표준이 데이터 교환 자동화를 위해 활용되어 왔으나, 이들은 주로 하위 수준의 데이터 전송에 초점을 맞추고 있으며, 설계 수준의 시뮬레이션 모델 구축을 위한 도메인별 어댑터 개발은 여전히 필수적이다~\cite{b11}.

최근 LLM의 코드 생성 능력을 활용하여 서로 다른 데이터 스키마 간 매핑 코드를 자동 합성하는 시도가 이루어지고 있다~\cite{b12}. 특히 ``산업 API를 위한 코드 생성(Code Generation for Industrial APIs)'' 연구는 자연어로 정의된 인터페이스 요구사항으로부터 실행 가능한 Python 스크립트를 생성하는 가능성을 입증하였다~\cite{b13}. 나아가, 시스템이 에러 로그를 자율적으로 분석하고 생성 코드를 수정하는 ``Self-healing'' 또는 ``Auto-Repair'' 메커니즘의 도입이 점차 중요해지고 있다~\cite{b14}. 이러한 자가 수리 접근법은 런타임 오류나 데이터 타입 불일치를 자동으로 해결하여 생성 코드의 실행 성공률을 크게 향상시킬 수 있다. 본 연구는 이러한 자가 수리 루프를 디지털 트윈 인터페이스에 적용하여, 기존 스크립트 연동부터 AI 기반 도구 합성까지 포괄하는 통합 연동 프레임워크의 강건성을 확보한다.

%% ============================================================
\section{제안 프레임워크}
\label{sec:framework}
%% ============================================================

본 장에서는 LLM을 활용하여 제조 라인 설계를 위한 구조화된 BOP 데이터를 자동 생성하고 분석 도구를 연동하는 프레임워크인 GDP의 아키텍처와 핵심 구성 요소를 제시한다.

\subsection{시스템 개요 및 사용자 인터페이스}

GDP는 Fig.~\ref{fig:overview} 상단에 도시된 바와 같이 세 개의 패널로 구성된 웹 기반 통합 환경이다.

\textbf{BOP 테이블(좌측 패널)}은 공정 목록, 사이클 타임, 배정된 설비·작업자·자재를 스프레드시트 형식으로 표시하며 각 셀을 직접 편집할 수 있다. \textbf{3D 레이아웃(중앙 패널)}은 공정 흐름과 리소스 배치를 3차원 공간에서 시각화한다. 각 공정은 박스로, 설비·작업자·자재는 타입별 마커로 렌더링되며, 공정 간 선후행 관계는 화살표로 표시된다. 사용자는 마우스 드래그로 모든 배치 객체에 대해 이동·회전·스케일링을 수행할 수 있다. \textbf{AI 어시스턴트(우측 패널)}은 자연어 채팅 인터페이스를 통해 BOP 생성·수정·질의를 LLM에 전달한다. LLM 프로바이더는 플러거블 방식으로 Gemini, GPT 등 다양한 모델을 지원한다.

세 패널은 양방향 동기화된다: 테이블에서 사이클 타임을 수정하면 3D 뷰가 즉시 갱신되고, 3D 뷰에서 공정을 드래그하면 테이블 좌표가 자동 갱신된다.

GDP는 단일 API에서 세 가지 상호작용 모드를 지원한다:
\begin{itemize}
    \item \textbf{BOP 생성}: 현재 BOP가 없을 때 사용자의 자연어 설명으로부터 새 BOP를 생성.
    \item \textbf{BOP 수정}: 현재 BOP가 존재할 때, 미수정 요소의 공간 좌표를 보존하면서 사용자 지시에 따라 수정.
    \item \textbf{질의응답}: BOP 데이터를 변경하지 않고 현재 BOP에 대한 분석적 질문(예: 병목 공정 식별)에 응답.
\end{itemize}

\subsection{공정 중심 데이터 모델}

Siemens가 제안한 BOP는 Plant $\rightarrow$ Line $\rightarrow$ Station $\rightarrow$ Operation의 심층 계층 구조를 갖지만, 이를 LLM 프롬프트에 그대로 포함하면 과도한 토큰 소비와 생성 정확도 저하를 초래한다. GDP는 이 문제를 해결하기 위해 Process를 핵심 단위로 하는 평탄화된 데이터 모델을 설계하였다.

BOP 데이터는 프로젝트 메타데이터(프로젝트명, 목표 UPH)와 함께 공정·설비·작업자·자재의 4종 마스터 리스트로 구성된다. 각 공정은 고유 ID, 공정명, 사이클 타임, 병렬 수, 선후행 관계를 보유하며, 리소스와의 연결은 별도의 배정(assignment) 엔티티를 통해 이루어진다. 배정 엔티티는 리소스 타입·참조 ID·수량과 함께 \textbf{공정 중심점 기준 상대 위치}를 보유한다. 리소스의 실제 3D 위치는 다음과 같이 계산된다:
\begin{equation}
\mathbf{p}_{\mathrm{actual}} = \mathbf{p}_{\mathrm{process}} + \mathbf{p}_{\mathrm{relative}}
\label{eq:coord}
\end{equation}

이 이중 좌표 체계를 통해 공정 단위를 이동할 때 하위 리소스 좌표를 개별 재계산할 필요가 없어 레이아웃 수정 비용이 크게 절감된다. 공정 간 선후행 관계는 임의 분기·합류를 허용하는 DAG 구조로 표현되며, DFS 기반 순환 감지가 실시간으로 적용된다.

\subsection{LLM 프로바이더 추상화}

모델 비의존적 BOP 생성을 위해 GDP는 프로바이더 추상화 계층을 도입한다. 추상 기본 클래스 \texttt{BaseLLMProvider}는 \texttt{generate()}, \texttt{generate\_json()}, \texttt{generate\_with\_retry()}의 핵심 메서드를 정의하며, Google Gemini(REST API)와 OpenAI(SDK 기반)에 대한 구체적 구현체가 이를 상속한다. 팩토리 함수가 모델 식별자에 따라 적절한 프로바이더를 동적 인스턴스화하여, 생성 파이프라인 수정 없이 LLM 백엔드를 전환할 수 있다.

\subsection{BOP 생성 파이프라인 (Step~1)}
\label{subsec:pipeline}

Fig.~\ref{fig:overview}의 Step~1은 자연어로부터 시뮬레이션 가능한 공정 설계를 생성하는 과정을 보여준다. BOP 생성 파이프라인은 다음 6단계로 구성된다:

\begin{enumerate}
    \item \textbf{프롬프트 구성}: 제조 공정 엔지니어 역할을 지시하는 도메인 특화 시스템 프롬프트를 사용자 요청과 결합한다. 시스템 프롬프트에는 JSON 출력 스키마, 공정 수 범위(3--6개), 리소스 배정 규칙(공정당 설비 1--3개, 작업자 1--2명, 자재 1--3종), 사이클 타임 범위(10--300초) 등 제조 도메인 제약 조건이 명시되어 있다.

    \item \textbf{LLM 추론}: 결정적 디코딩(temperature $= 0.0$)으로 선택된 프로바이더에 전송하며, JSON 파싱 실패 시 최대 3회 재시도를 지원한다.

    \item \textbf{수작업대 자동 보장}: 작업자가 배정되었으나 수작업대가 없는 공정에 \texttt{manual\_station} 타입 설비를 자동 삽입한다.

    \item \textbf{리소스 정렬}: 로봇 $\rightarrow$ 기계 $\rightarrow$ 수작업대 $\rightarrow$ 작업자 $\rightarrow$ 자재 순으로 결정적 정렬하여 일관된 3D 레이아웃을 보장한다.

    \item \textbf{DAG 기반 자동 레이아웃}: 위상 정렬 기반 레벨 할당으로 공정 좌표를 자동 계산한다.

    \item \textbf{3단계 검증}: Pydantic 스키마 적합성, 참조 무결성(모든 리소스 ID의 마스터 데이터 존재 여부), DFS 기반 DAG 순환 감지를 수행한다. 검증 실패 시 LLM에 재생성을 요청한다.
\end{enumerate}

\subsection{도구 연동 및 Auto-Repair (Step~2)}
\label{subsec:tools}

Fig.~\ref{fig:overview}의 Step~2는 생성된 BOP를 외부 분석 도구와의 연동을 통해 정제하는 과정을 보여준다. GDP는 전문가의 다양한 활용 시나리오를 지원하기 위해 세 가지 도구 연동 방식을 제공한다.

\textbf{방식 1: 기존 스크립트 어댑터 자동 생성.}
전문가가 이미 보유한 Python 분석 스크립트(argparse 기반 CLI 또는 JSON 파일 입출력 등)를 업로드하면, LLM이 소스 코드를 분석하여 입출력 스키마를 자동 추론한다. 이후 (i)~BOP 데이터를 도구가 기대하는 입력 형식으로 변환하는 전처리기와 (ii)~도구의 출력을 해석하여 BOP에 반영하는 후처리기, 두 개의 어댑터 함수를 LLM이 자동 생성한다.

\textbf{방식 2: AI 기반 분석 도구 로직 자체 합성.}
분석 스크립트가 없는 경우, 사용자가 자연어로 원하는 분석 내용을 설명하면(예: ``각 공정의 사이클 타임 대비 설비 가동률을 계산하라'') LLM이 분석 도구의 로직 자체를 Python 코드로 합성한다. GDP에 내장된 10종의 분석 도구(병목 분석, 라인 밸런싱, 에너지 추정, 레이아웃 최적화 등)가 이 방식으로 구현되었다.

\textbf{방식 3: 스키마 가이드 기반 고난도 최적화.}
고난도 최적화 도구(유전 알고리즘, 시뮬레이션 기반 최적화 등)를 연동할 때, GDP는 BOP 데이터 모델의 스키마 문서와 필드별 설명을 자동 제공하여, 사용자가 BOP의 내부 구조를 숙지하지 않아도 최적화 도구가 BOP 데이터에 접근할 수 있도록 추상화 계층을 제공한다.

\textbf{Auto-Repair 루프.}
세 가지 방식 모두에서, 어댑터 실행 시 런타임 오류가 발생하면 시스템이 오류 유형을 자동 분류(환경 제약 위반, 타입 불일치 등)하고, 어댑터 코드와 오류 진단 정보를 LLM에 전송하여 수정된 어댑터를 생성한다. 수정 코드로 재실행을 시도하며, 이 루프는 최대 $k$회 반복된다. 전체 과정은 타임스탬프, 오류 진단, 수리 이력과 함께 실행 로그에 기록되어 추적 가능성을 보장한다.

%% ============================================================
\section{실험}
\label{sec:experiments}
%% ============================================================

\subsection{실험 1: 제로샷 공정 생성 성능}

\subsubsection{실험 설정}

\textbf{정답 데이터.}
다양한 산업을 포괄하는 10종의 대표 제조 제품에 대한 정답 데이터셋을 공개 접근 가능한 기술 문헌~\cite{b15,b16,b17,b18,b19,b20,b21,b22,b23,b24}으로부터 직접 추출하여 구축하였다: EV 배터리 셀(P01, 11개 공정), 자동차 차체(P02, 4개), 스마트폰 SMT(P03, 7개), 반도체 후공정(P04, 7개), 태양광 PV 모듈(P05, 8개), 전기차 헤어핀 모터(P06, 9개), OLED 디스플레이(P07, 7개), 가정용 세탁기(P08, 6개), 연속식 제약 정제(P09, 6개), 타이어(P10, 6개). 총 71개 공정 단계로 구성되며, 한국어 및 영어로 주석되어 있다.

\textbf{평가 대상 모델.}
5개 LLM을 평가하였다:
\begin{enumerate}
    \item \textbf{Gemini 2.5 Flash} (Google, \texttt{gemini-2.5-flash})
    \item \textbf{GPT-5 Mini} (OpenAI, \texttt{gpt-5-mini})
    \item \textbf{o3-mini} (OpenAI, \texttt{o3-mini})
    \item \textbf{GPT-4o} (OpenAI, \texttt{gpt-4o})
    \item \textbf{GPT-4o Mini} (OpenAI, \texttt{gpt-4o-mini})
\end{enumerate}
모든 모델은 재현성을 위해 temperature $= 0.0$으로 호출하였으며, GDP의 실제 시스템 프롬프트를 사용하였다.

\textbf{평가 지표.}
두 가지 상호 보완적 지표를 사용한다:
\begin{itemize}
    \item \textbf{$n/M$ 정확도}: $M$을 정답 BOP 공정 수라 하자. 생성된 공정명은 한국어·영어 정답에 대한 최대 퍼지 유사도가 임계값 $\tau = 0.4$를 초과할 때 매칭으로 간주된다. 유사도 함수는 키워드 수준 Jaccard 유사도와 문자 수준 SequenceMatcher 비율을 결합한다.

    \item \textbf{시퀀스 매치}: 매칭된 $n$개 공정 중 올바르게 순서화된 쌍의 비율:
    \begin{equation}
        \text{SeqMatch} = \frac{\sum_{j < k} \mathbf{1}[i_j < i_k]}{\binom{n}{2}}
    \end{equation}
\end{itemize}

\subsubsection{결과}

Table~\ref{tab:model_averages}에 10개 제품에 대한 모델별 집계 결과를 제시한다.

\begin{table}[htbp]
\centering
\caption{10개 제품에 대한 모델별 평균 성능.}
\label{tab:model_averages}
\begin{tabular}{lcccc}
\hline
\textbf{모델} & \textbf{정확도(\%)} & \textbf{시퀀스(\%)} & \textbf{생성수} & \textbf{지연(s)} \\
\hline
Gemini 2.5 Flash & \textbf{81.2} & 75.8 & 9.1 & 31.4 \\
GPT-5 Mini       & 67.8 & 82.6 & 8.9 & 85.7 \\
o3-mini           & 51.8 & \textbf{83.3} & 4.7 & 29.3 \\
GPT-4o            & 40.3 & 73.3 & 4.2 & \textbf{12.1} \\
GPT-4o Mini       & 29.9 & 76.7 & 3.5 & 25.6 \\
\hline
\end{tabular}
\begin{tablenotes}
\small
\item 정확도: $n/M$ Accuracy; 시퀀스: Sequence Match; 생성수: 평균 생성 공정 수; 지연: 평균 응답 시간(초). 모든 모델의 성공률은 100\%.
\end{tablenotes}
\end{table}

Table~\ref{tab:product_results}에 각 모델의 제품별 정확도를 보고한다.

\begin{table}[htbp]
\centering
\caption{모델별 제품별 $n/M$ 정확도 (\%).}
\label{tab:product_results}
\begin{tabular}{lccccc}
\hline
\textbf{제품} & \textbf{Gemini} & \textbf{GPT-5m} & \textbf{o3m} & \textbf{4o} & \textbf{4om} \\
\hline
P01 (배터리)    & \textbf{100.0} & 90.9 & 45.5 & 27.3 & 36.4 \\
P02 (차체)      & 50.0  & \textbf{75.0} & 50.0 & 25.0 & 25.0 \\
P03 (SMT)       & \textbf{71.4} & \textbf{71.4} & 42.9 & 42.9 & 28.6 \\
P04 (반도체)    & \textbf{71.4} & 57.1 & 57.1 & 42.9 & 42.9 \\
P05 (태양광)    & \textbf{100.0} & 62.5 & 37.5 & 50.0 & 12.5 \\
P06 (모터)      & \textbf{66.7} & \textbf{66.7} & 44.4 & 22.2 & 44.4 \\
P07 (OLED)      & \textbf{85.7} & 71.4 & 57.1 & 42.9 & 42.9 \\
P08 (세탁기)    & \textbf{83.3} & 50.0 & 33.3 & 50.0 & 16.7 \\
P09 (제약)      & 83.3 & 33.3 & \textbf{100.0} & 50.0 & 16.7 \\
P10 (타이어)    & \textbf{100.0} & \textbf{100.0} & 50.0 & 50.0 & 33.3 \\
\hline
\textbf{평균}   & \textbf{81.2} & 67.8 & 51.8 & 40.3 & 29.9 \\
\hline
\end{tabular}
\begin{tablenotes}
\small
\item Gemini: Gemini 2.5 Flash; GPT-5m: GPT-5 Mini; o3m: o3-mini; 4o: GPT-4o; 4om: GPT-4o Mini.
\end{tablenotes}
\end{table}

\subsubsection{분석}

\textbf{공정 세분화 수준과 정확도.}
생성된 공정 수와 $n/M$ 정확도 사이에 강한 양의 상관관계가 관찰된다. Gemini 2.5 Flash와 GPT-5 Mini는 각각 평균 9.1개와 8.9개의 공정을 생성하여 정답 평균 7.1개에 근접한 세분화 수준을 달성하였다. 반면 GPT-4o, GPT-4o Mini, o3-mini는 평균 3.5--4.7개의 공정만을 생성하여, 핵심 제조 단계를 누락하는 지나치게 추상적인 표현을 산출하였다. 이는 제조 워크플로우를 세분화된 단계로 분해하는 모델의 능력이 BOP 생성 품질의 핵심 차별화 요소임을 시사한다.

\textbf{시퀀스 일관성.}
낮은 정확도에도 불구하고, o3-mini가 가장 높은 시퀀스 매치(83.3\%)를 달성하였으며, GPT-5 Mini(82.6\%)가 뒤를 이었다. 이는 추론 최적화 모델이 공정 포괄 범위가 제한적이더라도 올바른 공정 순서를 유지하는 경향이 있음을 나타낸다. 반대로, 가장 높은 정확도를 달성한 Gemini 2.5 Flash는 상대적으로 낮은 시퀀스 일관성(75.8\%)을 보여, 더 많은 공정 생성이 간혹 순서 오류를 유발할 수 있음을 시사한다.

\textbf{지연 시간--품질 트레이드오프.}
Gemini 2.5 Flash는 최고 정확도(81.2\%)를 적절한 지연 시간(31.4\,s)으로 달성하여, 실시간 인터랙티브 응용에 가장 적합한 백본으로 평가된다. GPT-4o는 가장 낮은 지연 시간(12.1\,s)을 보였으나 정확도는 4위에 그쳤다.

\textbf{도메인별 난이도.}
표준화되고 잘 문서화된 제조 공정을 가진 제품(P10: 타이어, P01: 배터리)에서 높은 정확도를 보인 반면, 전문적이거나 비표준적인 도메인(P02: 차체, P04: 반도체 후공정, P06: 헤어핀 스테이터)에서는 대부분의 모델이 60\% 미만의 점수를 기록하였다. 이는 특수 제조 도메인에 대한 검색 증강 생성(RAG)이나 Few-shot 프롬프팅의 필요성을 강조한다.

\subsection{전문가 인터뷰}

GDP의 도구 연동 기능에 대한 실용적 평가를 위해, 공정 최적화 분야의 박사급 전문가(SME) 1인을 대상으로 반구조화 인터뷰를 수행하였다. 인터뷰 대상자는 제조 시뮬레이션 및 공정 최적화 연구 경력 5년 이상의 전문가로, GDP의 세 가지 도구 연동 방식을 각각 시연받은 후 평가하였다.

\textbf{방식 1(기존 스크립트 어댑터 자동 생성)}에 대해, 전문가는 ``argparse 기반의 기존 최적화 스크립트를 BOP와 연동하기 위해 통상 2--3일 소요되던 어댑터 개발을 수 분 내에 완료할 수 있다''고 평가하였다. 특히 Auto-Repair가 타입 불일치와 같은 일반적 오류를 자동 해결하는 점이 실용적 가치가 높다고 언급하였다.

\textbf{방식 2(AI 기반 분석 도구 합성)}에 대해서는, ``탐색적 분석 단계에서 가설을 빠르게 검증하는 데 유용''하나, ``복잡한 수학적 모델이 필요한 경우 생성 코드의 정확성을 반드시 검증해야 한다''는 한계를 지적하였다.

\textbf{방식 3(스키마 가이드 기반 최적화)}에 대해서는, ``BOP 데이터 구조를 숙지하지 않은 연구자도 자신의 최적화 알고리즘을 GDP에 연동할 수 있는 진입 장벽 완화 효과''를 긍정적으로 평가하였다.

%% ============================================================
\section{결론}
\label{sec:conclusion}
%% ============================================================

본 연구는 LLM을 활용하여 제조 공정 설계와 도구 연동의 장벽을 제거하는 GDP 프레임워크를 제안하였다. GDP는 Siemens가 제안한 BOP 개념을 참조하되 공정(Process) 중심의 평탄화된 데이터 모델로 정제하여, LLM의 토큰 효율성과 생성 정확도를 동시에 확보하였다. DAG 기반 자동 레이아웃과 이중 좌표 체계로 수동 배치 비용을 최소화하고, 세 가지 도구 연동 방식(기존 스크립트 어댑터 자동 생성, AI 기반 도구 합성, 스키마 가이드 최적화)과 Auto-Repair 루프를 통해 이기종 분석 도구와의 연동 인터페이스 개발 부담을 대폭 경감하였다.

실험 결과, 5종의 LLM 중 Gemini 2.5 Flash가 10종 제조 제품에 대해 $n/M$ 정확도 81.2\%를 달성하여 제로샷 BOP 생성의 실현 가능성을 입증하였다. 전문가 인터뷰를 통해 도구 연동 기능의 실용적 가치를 확인하였으며, 특히 어댑터 자동 생성과 Auto-Repair의 조합이 인터페이스 개발 시간을 대폭 단축함을 검증하였다.

향후 연구에서는 (1)~Few-shot 프롬프팅 및 RAG를 통한 도메인 특화 정확도 향상, (2)~N:M 커버리지 매칭 도입을 통한 세분화 수준 불일치의 공정한 평가, (3)~출력 검증 결과를 피드백으로 활용하는 2차 수리 루프(validation-driven repair)를 통한 의미적 정확도 개선을 계획하고 있다. 또한, 비전문가가 GDP로 생성한 다양한 모의 생산 라인 데이터를 휴머노이드 로봇 및 피지컬 AI 시스템의 가상 학습 환경으로 활용하는 방안을 탐구할 예정이다.

GDP의 소스 코드와 실험 데이터는 \url{https://github.com/lee-geon-chang/gdp}에서 공개되어 있다.

%% ============================================================
%% 참고 문헌
%% ============================================================
\begin{thebibliography}{00}
\bibitem{b1} M.~Grieves and J.~Vickers, ``Digital twin: Mitigating unpredictable, undesirable emergent behavior in complex systems,'' in \textit{Transdisciplinary Perspectives on Complex Systems}, Springer, 2017, pp.~85--113.
\bibitem{b2} F.~Tao, J.~Cheng, Q.~Qi, M.~Zhang, H.~Zhang, and F.~Sui, ``Digital twin-driven product design, manufacturing and service with big data,'' \textit{Int.\ J.\ Adv.\ Manuf.\ Technol.}, vol.~94, pp.~3563--3576, 2018.
\bibitem{b3} ISO 23247-1:2021, ``Automation systems and integration --- Digital twin framework for manufacturing --- Part 1: Overview and general principles,'' 2021.
\bibitem{b4} T.~H.~J.~Uhlemann, C.~Lehmann, and R.~Steinhilper, ``The digital twin: Realizing the cyber-physical production system for Industry 4.0,'' \textit{Procedia CIRP}, vol.~61, pp.~335--340, 2017.
\bibitem{b5} F.~Ji, A.~Mardt, and T.~Nierhoff, ``Identifying inconsistencies in the design of large-scale casting systems --- An ontology-based approach,'' in \textit{Proc.\ IEEE Int.\ Conf.\ Autom.\ Sci.\ Eng.\ (CASE)}, 2022, pp.~1234--1240.
\bibitem{b6} T.~Brown \textit{et al.}, ``Language models are few-shot learners,'' in \textit{Proc.\ NeurIPS}, 2020, pp.~1877--1901.
\bibitem{b7} Siemens Digital Industries Software, ``Bill of Process (BOP) solution overview,'' 2023. [Online]. Available: \url{https://www.siemens.com}
\bibitem{b8} Y.~Sui, M.~Zhou, J.~Cai, S.~Hirber, and J.~Sun, ``Table meets LLM: Can large language models understand structured table data?'' in \textit{Proc.\ WSDM}, 2024.
\bibitem{b9} S.~Hegselmann, A.~Buber, and C.~Binnig, ``TabLLM: Few-shot classification of tabular data with large language models,'' \textit{arXiv:2210.10723}, 2022.
\bibitem{b10} G.~White, A.~Zink, L.~Codec\'{a}, and S.~Clarke, ``A digital twin smart city for citizen feedback,'' \textit{Cities}, vol.~110, p.~103064, 2021.
\bibitem{b11} AutomationML, ``IEC 62714: Engineering data exchange format for use in industrial automation systems engineering,'' 2018.
\bibitem{b12} M.~Chen \textit{et al.}, ``Evaluating large language models trained on code,'' \textit{arXiv:2107.03374}, 2021.
\bibitem{b13} N.~Vandemoortele, C.~Debruyne, and D.~O'Sullivan, ``Scalable table-to-knowledge graph matching from metadata using LLMs,'' in \textit{Proc.\ SemTab}, 2024.
\bibitem{b14} Y.~Brun \textit{et al.}, ``Engineering self-adaptive systems through feedback loops,'' in \textit{Software Engineering for Self-Adaptive Systems}, Springer, 2013, pp.~48--70.
\bibitem{b15} VDMA, ``Production process of a lithium-ion battery cell,'' [Online]. Available: \url{https://vdma-industryguide.com}
\bibitem{b16} Toyota Motor Corp., ``Plant tours: Stamping,'' [Online]. Available: \url{https://global.toyota/en/company/plant-tours/stamping/}
\bibitem{b17} PCBOnline, ``SMT manufacturing process,'' [Online]. Available: \url{https://www.pcbonline.com/blog/smt-manufacturing-process.html}
\bibitem{b18} Amkor Technology, ``System-in-Package technology,'' 2019.
\bibitem{b19} NREL, ``PV module manufacturing process flow,'' NREL/TP-520-32943, 2003.
\bibitem{b20} thyssenkrupp, ``Hairpin stator manufacture,'' Tech Days, 2022.
\bibitem{b21} UBI Research, ``AMOLED manufacturing process report,'' 2021.
\bibitem{b22} I.~Khan \textit{et al.}, ``Design for manufacturing and assembly of washing machine,'' \textit{ResearchGate}, 2017.
\bibitem{b23} QbD Works, ``Continuous manufacturing --- FDA perspective,'' 2014.
\bibitem{b24} Nexen Tire Corp., ``Tire manufacturing process,'' [Online]. Available: \url{https://www.nexentire.com}
\end{thebibliography}
